\documentclass[../../main.tex]{subfiles}
\begin{document}

{\bf [解]}
$x^4$ を $p(x)$ でわった商を $A(x)$，あまりを $Q(x)$ とする．また $x^5$ を $p(x)$ でわった商を $B(x)$ とする．
すると題意より
\begin{align}
\begin{cases}
x^4=A(x)\cdot p(x)+Q(x) \\
x^5=B(x)\cdot p(x)+Q(x)
\end{cases} \label{eq:1}
\end{align}
と書ける．\cref{eq:1}の辺々引いて，
\begin{align}
x^4(x-1)=\{B(x)-A(x)\}p(x)
\end{align}
だから $x^4(x-1)$ は $p(x)$ で割り切れるので，$a$ を係数として3次式$p(x)$は
\begin{align}
p(x)=ax^3, \quad a(x-1) x^2
\end{align}
と書ける．ところが$p(x)=ax^3$の時は\cref{eq:1}より$Q(x)=0$となり，これは題意に反して不適である．
よって
\begin{align}
 p(x)=a(x-1) x^2 \label{eq:2}
\end{align}
である．


次に $f(x)$ を $p(x)$ でわった商 $C(x)$，$f(x)\cdot x$ を $p(x)$ でわった商を $D(x)$ とする．
題意より
\begin{align}
\begin{cases}
f(x)=C(x)\cdot p(x)+r(x) & \cdots \text{②} \\
xf(x)=D(x)\cdot p(x) & \cdots \text{③}
\end{cases} \label{eq:3}
\end{align}
と書ける．\cref{eq:3}の第一式に$x$をかけたものと第二式を辺々引いて，
\begin{align}
0=\{x\cdot C(x)-D(x)\}\cdot p(x)+xr(x)
\end{align}
だから $xr(x)$ は $p(x)$ で割り切れ，かつ$p(x)$が3次だから $r(x)$ は2次以下である．
よって\cref{2}および題意から最高次の係数が$1$であることから
\begin{align}
r(x)=x(x-1)
\end{align}
である．

\newpage

\end{document}
